% Part: first-order-logic
% Chapter: models-theories
% Section: theories

\documentclass[../../../include/open-logic-section]{subfiles}

\begin{document}

\olfileid{fol}{mat}{the}

\olsection{نمونه‌هایی از نظریه‌های مرتبهٔ اول}

\begin{ex}
نظریهٔ ترتیب‌های خطی اکید در زبان~$\Lang L_<$ به‌وسیلهٔ مجموعهٔ زیر
اصل‌موضوع‌بندی می‌شود:
\begin{align*}
\{\quad & \lforall[x][\lnot x < x], \\
& \lforall[x][\lforall[y][((x < y \lor y <
    x) \lor x = y)]], \\
& \lforall[x][\lforall[y][\lforall[z][((x < y
      \land y < z) \lif x < z)]]] \quad \}
\end{align*}
این نظریه !!{structure}s مورد نظر را به‌طور کامل مشخصه‌سازی می‌کند: هر
ترتیب خطی اکید مدلی از این دستگاه اصول موضوع است، و برعکس، اگر $R$
ترتیبی خطی روی مجموعهٔ $X$ باشد، آنگاه ساختار~$\Struct M$ با
$\Domain M = X$ و $\Assign{<}{M} = R$ مدلی از این نظریه است.
\end{ex}

\begin{ex}
نظریهٔ گروه‌ها در زبان شامل~$\Obj 1$ (!!{constant}) و~$\cdot$
(!!{function} دوموضعی) به‌وسیلهٔ اصول موضوع زیر اصل‌موضوع‌بندی می‌شود:
\begin{align*}
& \lforall[x][\eq[(x \cdot \Obj 1)][x]]\\
& \lforall[x][\lforall[y][\lforall[z][\eq[(x \cdot (y \cdot z))][((x
          \cdot y) \cdot z)]]]]\\
& \lforall[x][\lexists[y][\eq[(x \cdot y)][\Obj 1]]]
\end{align*}
\end{ex}

\begin{ex}
نظریهٔ حساب پئانو به‌وسیلهٔ جمله‌های زیر در زبان
حساب~$\Lang L_A$ اصل‌موضوع‌بندی می‌شود:
\begin{align*}
& \lforall[x][\lforall[y][(\eq[x'][y'] \lif \eq[x][y])]]\\
& \lforall[x][\eq/[\Obj 0][x']]\\
& \lforall[x][\eq[(x + \Obj 0)][x]]\\
& \lforall[x][\lforall[y][\eq[(x + y')][(x + y)']]]\\
& \lforall[x][\eq[(x \times \Obj 0)][\Obj 0]]\\
& \lforall[x][\lforall[y][\eq[(x \times y')][((x \times y) + x)]]]\\
& \lforall[x][\lforall[y][(x < y \liff \lexists[z][\eq[(z' + x)][y])]]]\\
\intertext{به‌علاوهٔ همهٔ جمله‌هایی به صورت}
& (!A(\Obj 0) \land \lforall[x][(!A(x) \lif !A(x'))]) \lif \lforall[x][!A(x)]
\end{align*}
چون بی‌نهایت جمله به صورت اخیر وجود دارد، این دستگاه اصول موضوع
نامتناهی است. صورت اخیر \emph{طرح‌وارهٔ استقرا} نامیده می‌شود. (در واقع،
طرح‌وارهٔ استقرا اندکی پیچیده‌تر از آن است که اینجا وانمود کرده‌ایم.)

اصل موضوع آخر یک \emph{تعریف صریح} از~$<$ است.
\end{ex}

\begin{ex}
نظریهٔ مجموعه‌های محض در مبانی ریاضیات (و در فلسفهٔ ریاضیات) نقشی مهم
دارد. مجموعه‌ای محض است اگر همهٔ !!{element}s آن نیز مجموعه‌های محض
باشند. بنابراین مجموعهٔ تهی محض به شمار می‌آید، اما مجموعه‌ای که
!!a{element} دارد که مجموعه نیست، محض نخواهد بود. پس مجموعه‌های محض
همان‌هایی هستند که تنها از مجموعهٔ تهی و بدون هیچ «فرد اولیه‌ای»، یعنی
بدون اشیایی که خود مجموعه نیستند، ساخته شده‌اند.

مجموعهٔ زیر را می‌توان دستگاهی از اصول موضوع برای نظریه‌ای دربارهٔ
مجموعه‌های محض دانست:
\begin{align*}
& \lexists[x][\lnot \lexists[y][y \in x]]\\
& \lforall[x][\lforall[y][(\lforall[z][(z \in x \liff z \in y)] \lif
\eq[x][y])]]\\
& \lforall[x][\lforall[y][\lexists[z][\lforall[u][(u \in z \liff
(\eq[u][x] \lor \eq[u][y]))]]]]\\
& \lforall[x][\lexists[y][\lforall[z][(z \in y \liff \lexists[u][(z \in
        u \land u \in x)])]]]\\
\intertext{به‌علاوهٔ همهٔ جمله‌هایی به صورت} &
\lexists[x][\lforall[y][(y \in x \liff !A(y))]]
\end{align*}
اصل موضوع نخست می‌گوید مجموعه‌ای بدون !!{element} وجود دارد (یعنی
$\emptyset$ وجود دارد)؛ اصل دوم می‌گوید مجموعه‌ها گسترشی‌اند؛ اصل سوم
می‌گوید برای هر دو مجموعهٔ $X$ و~$Y$، مجموعهٔ~$\{X, Y\}$ وجود دارد؛ و
اصل چهارم می‌گوید برای هر مجموعهٔ $X$، مجموعهٔ~$\cup X$ وجود دارد که در
آن $\cup X$ اجتماع همهٔ اعضای~$X$ است.

!!{sentence}sی که در پایان ذکر شدند، در مجموع \emph{طرح‌وارهٔ
مجموعه‌سازی ساده‌انگارانه} نامیده می‌شوند. این طرح‌واره در اصل می‌گوید
برای هر~$!A(x)$، مجموعهٔ~$\Setabs{x}{!A(x)}$ وجود دارد---و ازاین‌رو در
نگاه نخست اصلی صادق، سودمند و شاید حتی ضروری به نظر می‌رسد. آن را
«ساده‌انگارانه» می‌نامند، زیرا چنان‌که معلوم می‌شود، این نظریه را
ارضا‌ناپذیر می‌سازد: اگر $!A(y)$ را~$\lnot y \in y$ بگیریم، به
!!{sentence} زیر می‌رسیم:
\[
\lexists[x][\lforall[y][(y \in x \liff \lnot y \in y)]]
\]
و این !!{sentence} در هیچ !!{structure}ای ارضا نمی‌شود.
\end{ex}

\begin{ex}
در حوزهٔ \emph{جزءشناسی}، رابطهٔ \emph{جزءبودن} رابطه‌ای بنیادی است.
همانند نظریه‌های مجموعه‌ها، نظریه‌هایی نیز دربارهٔ جزءبودن وجود دارند
که برداشت‌های گوناگون (و گاه متعارض) از این رابطه را اصل‌موضوع‌بندی
می‌کنند.

زبان جزءشناسی تنها یک نماد محمول دوموضعی~$\Obj P$ دارد و
$\Atom{\Obj P}{x, y}$ «به این معناست» که $x$ جزئی از~$y$ است. اگر چنین
تعبیری را در نظر داشته باشیم، !!a{structure} برای این زبان
\emph{ساختار جزءبودن} نامیده می‌شود. البته هر ساختاری برای یک
محمول دوموضعی واقعاً شایستهٔ این نام نیست. برای آنکه بختی برای
مشخصه‌سازی «جزءبودن» داشته باشیم، $\Assign{\Obj P}{M}$ باید شرایطی را
برآورده کند که می‌توانیم آن‌ها را به‌عنوان اصول موضوع نظریه‌ای دربارهٔ
جزءبودن وضع کنیم. برای نمونه، جزءبودن یک ترتیب جزئی روی اشیاست: هر شیء
جزئی (ولو جزئی \emph{ناسره}) از خودش است؛ دو شیء متفاوت نمی‌توانند جزء
یکدیگر باشند؛ و جزءِ جزءِ یک شیء، خود جزئی از آن شیء است. توجه کنید که
«جزئی از ... است» به این معنا به «زیرمجموعهٔ ... است» شباهت دارد، اما به
«عضوِ ... است»، که نه بازتابی است و نه تراگذری، شباهت ندارد.
\begin{align*}
& \lforall[x][\Atom{\Obj P}{x,x}] \\
& \lforall[x][\lforall[y][((\Part{x}{y} \land \Part{y}{x})
      \lif \eq[x][y])]] \\
& \lforall[x][\lforall[y][\lforall[z][((\Part{x}{y} \land
        \Part{y}{z}) \lif \Part{x}{z})]]]\\
\intertext{افزون بر این، هر دو شیء یک مجموع جزءشناختی دارند (شیئی که
  این دو شیء را جزء خود دارد و از این حیث کمینه است).}  &
\lforall[x][\lforall[y][\lexists[z][\lforall[u][(\Part{z}{u} \liff
        (\Part{x}{u} \land \Part{y}{u}))]]]]
\end{align*}
این‌ها تنها برخی از اصول بنیادی جزءبودن‌اند که متافیزیک‌دانان در نظر
گرفته‌اند. بااین‌حال، صورت‌بندی یا نوشتن اصول بیشتر بی‌آنکه نخست چند
رابطهٔ تعریف‌شده معرفی کنیم، به‌سرعت دشوار می‌شود. برای نمونه، بیشتر
متافیزیک‌دانانی که به جزءشناسی علاقه‌مندند اصل زیر را نیز معتبر
می‌دانند: هرگاه شیء~$x$ جزء سره‌ای چون~$y$ داشته باشد، جزء دیگری
چون~$z$ نیز دارد که هیچ جزء مشترکی با~$y$ ندارد، و هم‌جوشی $y$ و~$z$
برابر با~$x$ است.
\end{ex}

\end{document}
